\chapter{Error Handling}
\label{ch:errors}

\haira{} uses explicit error handling inspired by Go. Errors are values, not exceptions. Functions that can fail return both a result and an optional error.

\section{Philosophy}

\begin{itemize}
    \item \textbf{Errors are values} -- Not special control flow
    \item \textbf{Explicit handling} -- Callers must handle or propagate errors
    \item \textbf{No hidden control flow} -- No exceptions or stack unwinding
    \item \textbf{Composable} -- Errors can be wrapped and chained
\end{itemize}

\section{The Error Type}

The built-in \type{Error} type represents errors:

\begin{lstlisting}
struct Error {
    message: string
    code: int?
    source: Error?
}
\end{lstlisting}

\subsection{Creating Errors}

\begin{lstlisting}
// Simple error
err = Error{message: "something went wrong"}

// With error code
err = Error{message: "permission denied", code: 403}

// Wrapped error (error chaining)
err = Error{
    message: "failed to read config",
    source: original_error
}
\end{lstlisting}

\section{Error Return Pattern}

Functions that can fail return a tuple of \code{(result, Error?)}:

\begin{lstlisting}
fn divide(a: int, b: int) -> (int, Error?) {
    if b == 0 {
        return 0, Error{message: "division by zero"}
    }
    return a / b, nil
}

fn read_file(path: string) -> (string, Error?) {
    // Implementation
}

fn parse_int(s: string) -> (int, Error?) {
    // Implementation
}
\end{lstlisting}

\section{Handling Errors}

\subsection{Basic Error Checking}

\begin{lstlisting}
result, err = divide(10, 0)
if err != nil {
    io.println("Error: " + err.message)
    return
}
io.println("Result: " + to_string(result))
\end{lstlisting}

\subsection{Propagating Errors}

\begin{lstlisting}
fn process_file(path: string) -> (Data, Error?) {
    content, err = read_file(path)
    if err != nil {
        return Data{}, err  // Propagate error
    }

    data, err = parse_data(content)
    if err != nil {
        return Data{}, err  // Propagate error
    }

    return data, nil
}
\end{lstlisting}

\subsection{Wrapping Errors}

Add context when propagating:

\begin{lstlisting}
fn load_config(path: string) -> (Config, Error?) {
    content, err = read_file(path)
    if err != nil {
        return Config{}, Error{
            message: "failed to load config from " + path,
            source: err
        }
    }

    config, err = parse_config(content)
    if err != nil {
        return Config{}, Error{
            message: "invalid config format",
            source: err
        }
    }

    return config, nil
}
\end{lstlisting}

\section{Error Checking Patterns}

\subsection{Guard Pattern}

Check errors early and return:

\begin{lstlisting}
fn process(input: string) -> (Output, Error?) {
    // Guard: validate input
    if string.len(input) == 0 {
        return Output{}, Error{message: "input cannot be empty"}
    }

    // Guard: check prerequisites
    ready, err = check_ready()
    if err != nil {
        return Output{}, err
    }
    if not ready {
        return Output{}, Error{message: "system not ready"}
    }

    // Main logic
    result = do_processing(input)
    return result, nil
}
\end{lstlisting}

\subsection{Defer for Cleanup}

Use \keyword{defer} to ensure cleanup happens even on error:

\begin{lstlisting}
fn with_file(path: string) -> (string, Error?) {
    file, err = fs.open(path)
    if err != nil {
        return "", err
    }
    defer fs.close(file)  // Always closes, even on error

    content, err = fs.read_all(file)
    if err != nil {
        return "", err  // file is closed by defer
    }

    return content, nil  // file is closed by defer
}
\end{lstlisting}

\section{Error Utilities}

\subsection{Checking Error Types}

\begin{lstlisting}
fn is_not_found(err: Error?) -> bool {
    if err == nil {
        return false
    }
    return err.code == 404 or string.contains(err.message, "not found")
}

fn is_permission_denied(err: Error?) -> bool {
    if err == nil {
        return false
    }
    return err.code == 403
}
\end{lstlisting}

\subsection{Error Chain Traversal}

\begin{lstlisting}
fn root_cause(err: Error?) -> Error? {
    if err == nil {
        return nil
    }
    current = err
    while current.source != nil {
        current = current.source
    }
    return current
}

fn error_chain(err: Error?) -> []Error {
    errors = []
    current = err
    while current != nil {
        errors = array.push(errors, current)
        current = current.source
    }
    return errors
}
\end{lstlisting}

\section{Option Types and Errors}

For operations that may not find a value (not an error condition), use option types:

\begin{lstlisting}
// Use option when "not found" is normal
fn find_user(id: int) -> User? {
    // Returns nil if not found
}

// Use error when "not found" is exceptional
fn get_user(id: int) -> (User, Error?) {
    // Returns error if not found
}
\end{lstlisting}

\subsection{Converting Between Options and Errors}

\begin{lstlisting}
// Option to error
fn require_user(id: int) -> (User, Error?) {
    user = find_user(id)
    if user == nil {
        return User{}, Error{message: "user not found", code: 404}
    }
    return user, nil
}

// Error to option (ignoring error)
fn try_parse_int(s: string) -> int? {
    result, err = parse_int(s)
    if err != nil {
        return nil
    }
    return result
}
\end{lstlisting}

\section{Panic and Recovery}

For truly unrecoverable situations, use \code{panic}:

\begin{lstlisting}
fn assert(condition: bool, message: string) {
    if not condition {
        panic(message)
    }
}

fn unwrap<T>(opt: T?, message: string) -> T {
    if opt == nil {
        panic(message)
    }
    return opt
}
\end{lstlisting}

\begin{warning}
\code{panic} terminates the program. Use it only for programming errors, not for expected failure conditions. Prefer returning errors for all recoverable situations.
\end{warning}

\subsection{When to Use Panic}

\begin{itemize}
    \item Programming bugs (assertions)
    \item Invariant violations
    \item Unrecoverable initialization failures
    \item Unreachable code paths
\end{itemize}

\subsection{When to Use Errors}

\begin{itemize}
    \item File not found
    \item Network failures
    \item Invalid user input
    \item Resource exhaustion
    \item Any expected failure condition
\end{itemize}

\section{Best Practices}

\subsection{Return Errors, Don't Log and Continue}

\begin{lstlisting}
// Bad: logs and continues with invalid state
fn process(data: Data) -> Result {
    result, err = validate(data)
    if err != nil {
        log.error("validation failed: " + err.message)
        // Continues with unvalidated data!
    }
    // ...
}

// Good: returns error to caller
fn process(data: Data) -> (Result, Error?) {
    result, err = validate(data)
    if err != nil {
        return Result{}, Error{
            message: "validation failed",
            source: err
        }
    }
    // ...
}
\end{lstlisting}

\subsection{Add Context When Wrapping}

\begin{lstlisting}
// Bad: just propagates
if err != nil {
    return Result{}, err
}

// Good: adds context
if err != nil {
    return Result{}, Error{
        message: "failed to process user " + to_string(user_id),
        source: err
    }
}
\end{lstlisting}

\subsection{Handle Errors at the Right Level}

\begin{lstlisting}
// Low-level: specific error
fn read_config_file(path: string) -> ([]byte, Error?) {
    // Returns "file not found" or "permission denied"
}

// Mid-level: wrapped with context
fn load_config(path: string) -> (Config, Error?) {
    data, err = read_config_file(path)
    if err != nil {
        return Config{}, Error{
            message: "could not load config",
            source: err
        }
    }
    // Parse...
}

// High-level: user-friendly handling
fn main() {
    config, err = load_config("config.json")
    if err != nil {
        io.println("Error: " + err.message)
        // Could show root_cause for debugging
        os.exit(1)
    }
    // Continue...
}
\end{lstlisting}

\section{Error Handling Grammar}

Error handling uses standard language constructs:

\begin{lstlisting}[language=EBNF]
error_return = "(" type "," "Error?" ")" ;

error_check = "if" identifier "!=" "nil" block ;

error_create = "Error" "{"
               "message" ":" expression
               [ "," "code" ":" expression ]
               [ "," "source" ":" expression ]
               "}" ;
\end{lstlisting}
